\chapter{Fecal Occult Blood Test (FOBT)}
\section*{What Is This Test?} FOBT detects small amounts of blood in stool that cannot be seen. It is used mainly for colorectal-cancer screening.
\section*{Alternative Names} Stool occult blood test; guaiac fecal occult blood test; gFOBT; fecal immunochemical test (FIT).
\section*{Principle} Guaiac tests detect heme chemically; FIT uses antibodies against human hemoglobin and is more specific for lower gastrointestinal bleeding.
\section*{Why This Test Is Ordered} It screens average-risk adults for colorectal cancer and may investigate suspected gastrointestinal bleeding, depending on the clinical setting.
\section*{Primary vs Secondary Test} Stool screening is primary for selected screening programs; a positive result requires colon evaluation, not a cancer diagnosis.
\section*{How This Test Is Performed} The patient collects one or more stool samples according to kit instructions and returns them promptly.
\section*{What Preparation Is Needed} FIT usually has few dietary restrictions; gFOBT may require diet and medicine instructions. Follow the kit exactly.
\section*{Understanding Results} A positive result can reflect cancer, polyps, hemorrhoids, ulcer, or other bleeding. A negative result does not eliminate cancer risk.
\section*{Follow-up Tests} Positive screening requires diagnostic colonoscopy; repeat testing is not a substitute. CBC and other studies may assess anemia or bleeding.
\section*{References} \begin{enumerate}\item MedlinePlus. Fecal Occult Blood Test.\item U.S. Preventive Services Task Force. Colorectal cancer screening recommendations.\end{enumerate}
\clearpage\section*{Understanding Results and Follow-up}
The laboratory report should identify the specimen, method, units, reference interval or decision limit, and whether the result is positive, negative, abnormal, or inconclusive. Reference intervals vary with age, sex, pregnancy, specimen type, assay, and laboratory; a result outside a range is not automatically a diagnosis, and a result within a range does not always exclude disease. Discuss unexpected results with the ordering clinician, who can decide whether repeat or confirmatory testing, treatment, or referral is appropriate. Do not change medicines or delay urgent care based on an isolated result.
\section*{Regulatory Status and Authorization Date}This chapter describes a test category, not a specific commercial product. FDA, CDSCO, EU IVDR, and MHRA approval or registration dates therefore cannot be assigned without the assay manufacturer, product name, instrument, and intended use. For laboratory-developed tests, an individual agency approval date may not exist. See the Preface for the interpretation of regulatory status.
\clearpage
